\section{The full-system gap on two replicas}
\label{sec:doubled_system_gap}

Let $\mathcal H$ have a finite orthonormal basis indexed by $n$.
For $H\in\mathcal L(\mathcal H)$, $\Omega\in\mathcal H$, and
$E_0,\Delta\in\R$, write
\begin{align}
    P&=\ketbra{\Omega}{\Omega}, & Q&=I-P, &
    R&=H-E_0I-\Delta Q.
    \label{eq:doubled_gap_residual}
\end{align}
The lower-gap assertion for the two replicas in
\cite{OpenAI2026PolynomialPEPS} follows by decomposing their corresponding
residual into positive semidefinite terms.
% Source: pinned 02-information.tex, lines 415--425, the unnumbered passage
% immediately after eq:info-reset-overlap. The normalized eigenvector,
% ground-energy, lower-gap, and strictly positive-gap uniqueness assertions
% are covered here. Interaction support and exact gap equality are not claimed.

\begin{definition}[Doubled vector and Hamiltonian]
    \label{def:doubled_system}
    \lean{Matrix.doubledVector, Matrix.doubledHamiltonian}
    \leanok
    On the original doubled space $\mathcal H\otimes\mathcal H$, define
    $\Psi=\Omega\otimes\Omega$ and
    $\mathbb H=H\otimes I+I\otimes H$. In coordinates,
    $\Psi_{(i,j)}=\Omega_i\Omega_j$.
\end{definition}

\begin{theorem}[Product projector, normalization, and energy]
    \label{thm:doubled_vector_properties}
    \lean{Matrix.vecMulVec_doubledVector, Matrix.sum_normSq_doubledVector,
        Matrix.norm_doubledVector, Matrix.doubledHamiltonian_mulVec}
    \leanok
    \uses{def:doubled_system}
    The doubled vector satisfies
    \begin{align}
        \ketbra{\Psi}{\Psi}&=P\otimes P, &
        \sum_{i,j}|\Psi_{(i,j)}|^2
            &=\left(\sum_i|\Omega_i|^2\right)^2.
        \label{eq:doubled_vector_product}
    \end{align}
    Consequently, $\|\Omega\|=1$ implies $\|\Psi\|=1$.
    If $H\Omega=E_0\Omega$, then $\mathbb H\Psi=2E_0\Psi$.
\end{theorem}
\begin{proof}\leanok
    \uses{def:doubled_system}
    The projector entries are
    $\Omega_i\Omega_j\overline{\Omega_k\Omega_l}
    =P_{ik}P_{jl}$, and the squared norm factors as a product of two
    finite sums. For the eigenvector assertion,
    \begin{align}
        \mathbb H\Psi
            &=(H\Omega)\otimes\Omega+\Omega\otimes(H\Omega)
              =2E_0\Psi.
    \end{align}
\end{proof}

\begin{theorem}[Preservation of the nonnegative lower gap]
    \label{thm:doubled_lower_gap}
    \lean{Matrix.doubledHamiltonian_gap_eq, Matrix.PosSemidef.doubled_gap}
    \leanok
    \uses{def:doubled_system}
    With the notation of~(\ref{eq:doubled_gap_residual}), the algebraic identity
    \begin{align}
        \mathbb H-2E_0I-\Delta(I-\ketbra{\Psi}{\Psi})
            &=R\otimes I+I\otimes R+\Delta(Q\otimes Q)
        \label{eq:doubled_gap_decomposition}
    \end{align}
    holds for every $H,\Omega,E_0,\Delta$.
    If $R\ge0$, $\|\Omega\|=1$, and $\Delta\ge0$, then
    \begin{align}
        \mathbb H-2E_0I&\ge\Delta(I-\ketbra{\Psi}{\Psi}).
        \label{eq:doubled_gap_bound}
    \end{align}
    No eigenvector assumption is needed for this inequality.
\end{theorem}
\begin{proof}\leanok
    \uses{thm:doubled_vector_properties}
    Substitute $\ketbra{\Psi}{\Psi}=P\otimes P$ and expand both sides
    of~(\ref{eq:doubled_gap_decomposition}). Normalization gives $Q\ge0$:
    for every $v\in\mathcal H$,
    $\langle v,Qv\rangle=\|v\|^2-|\langle\Omega,v\rangle|^2\ge0$.
    Tensor products and sums of positive semidefinite operators are
    positive semidefinite. Thus each term on the right of
    (\ref{eq:doubled_gap_decomposition}) is positive semidefinite when
    $R\ge0$ and $\Delta\ge0$, proving~(\ref{eq:doubled_gap_bound}).
\end{proof}

\begin{theorem}[Unitary transport of a lower gap]
    \label{thm:doubled_unitary_transport}
    \lean{Matrix.norm_toLp_mulVec_of_unitary, Matrix.gap_unitary_conj_eq,
        Matrix.PosSemidef.gap_unitary_conj,
        Matrix.mulVec_unitary_conj_eigenvector}
    \leanok
    Let $U$ be unitary on $\mathcal H$. Then $\|U\Omega\|=\|\Omega\|$ and
    \begin{align}
        UHU^\dagger-E_0I-\Delta(I-\ketbra{U\Omega}{U\Omega})
            &=URU^\dagger.
        \label{eq:doubled_gap_conjugation}
    \end{align}
    Hence $R\ge0$ implies that the left side of
    (\ref{eq:doubled_gap_conjugation}) is positive semidefinite.
    If $H\Omega=E_0\Omega$, then
    $(UHU^\dagger)(U\Omega)=E_0(U\Omega)$.
    These statements hold for every real $E_0,\Delta$, without a
    normalization assumption.
\end{theorem}
\begin{proof}\leanok
    The identities $U^\dagger U=UU^\dagger=I$ give
    $\langle U\Omega,U\Omega\rangle=\langle\Omega,\Omega\rangle$ and
    $\ketbra{U\Omega}{U\Omega}=UPU^\dagger$.
    Expanding $URU^\dagger$ proves~(\ref{eq:doubled_gap_conjugation});
    congruence by $U$ preserves positive semidefiniteness. Finally,
    $(UHU^\dagger)(U\Omega)=UH\Omega=E_0(U\Omega)$.
\end{proof}

\begin{theorem}[Uniqueness from a strictly positive lower gap]
    \label{thm:positive_gap_uniqueness}
    \lean{Matrix.PosSemidef.eq_inner_smul_of_gap,
        Matrix.PosSemidef.eigenspace_eq_span_of_gap,
        Matrix.PosSemidef.finrank_eigenspace_eq_one_of_gap}
    \leanok
    Suppose $R\ge0$, $\Delta>0$, and $\|\Omega\|=1$.
    Every vector $\psi$ satisfying $H\psi=E_0\psi$ obeys
    $\psi=\langle\Omega,\psi\rangle\Omega$.
    The vector $\psi$ may be zero or unnormalized; no eigenvector equation
    on $\Omega$ is required for this projection identity.
    If also $H\Omega=E_0\Omega$, then
    \begin{align}
        \ker(H-E_0I)&=\spn_{\C}\{\Omega\}, &
        \dim_{\C}\ker(H-E_0I)&=1.
        \label{eq:positive_gap_ground_line}
    \end{align}
\end{theorem}
\begin{proof}\leanok
    For $H\psi=E_0\psi$, the lower-gap inequality gives
    \begin{align}
        \Delta\bigl(\|\psi\|^2-|\langle\Omega,\psi\rangle|^2\bigr)
            &\le \operatorname{Re}\langle\psi,H\psi\rangle
                -E_0\|\psi\|^2=0.
        \label{eq:positive_gap_energy}
    \end{align}
    Since $\Delta>0$ and $\|\Omega\|=1$,
    \begin{align}
        \|\psi-\langle\Omega,\psi\rangle\Omega\|^2
            &=\|\psi\|^2-|\langle\Omega,\psi\rangle|^2\le0.
        \label{eq:positive_gap_projection}
    \end{align}
    Thus the residual vector in~(\ref{eq:positive_gap_projection}) is zero.
    Every vector in $\ker(H-E_0I)$ consequently belongs to
    $\spn_{\C}\{\Omega\}$. The reverse inclusion follows from
    $H\Omega=E_0\Omega$. Finally, $\|\Omega\|=1$ implies $\Omega\ne0$,
    so its span has dimension one, proving~(\ref{eq:positive_gap_ground_line}).
\end{proof}

\begin{theorem}[The physical partial swap preserves the doubled gap]
    \label{thm:doubled_partial_swap_gap}
    \lean{Matrix.doubled_gap_partialSwap}
    \leanok
    \uses{def:doubled_system, def:physical_buffer_compression}
    Suppose $\mathcal H=\mathcal H_A\otimes\mathcal H_B$,
    $\|\Omega\|=1$, $H\Omega=E_0\Omega$, $\Delta\ge0$, and
    $H-E_0I\ge\Delta(I-\ketbra{\Omega}{\Omega})$.
    Let $F=F_A$ be the physical partial swap from
    Definition~\ref{def:physical_buffer_compression}, and put
    $\Phi=F\Psi$ and $\mathbb H'=F\mathbb HF^\dagger$. Then
    \begin{align}
        \|\Phi\|&=1, & \mathbb H'\Phi&=2E_0\Phi,
        \label{eq:doubled_swap_conclusions}\\
        \mathbb H'-2E_0I&\ge\Delta(I-\ketbra{\Phi}{\Phi}). \notag
    \end{align}
    Both vectors and both Hamiltonians act on the same doubled physical
    space $(\mathcal H_A\otimes\mathcal H_B)^{\otimes2}$.
\end{theorem}
\begin{proof}\leanok
    \uses{thm:physical_buffer_contraction, thm:doubled_vector_properties,
        thm:doubled_lower_gap, thm:doubled_unitary_transport}
    The partial swap is unitary by
    Theorem~\ref{thm:physical_buffer_contraction}.
    Apply Theorem~\ref{thm:doubled_vector_properties} to obtain the
    normalized vector $\Psi$ and its eigenvalue $2E_0$.
    Theorem~\ref{thm:doubled_lower_gap} gives
    (\ref{eq:doubled_gap_bound}), and
    Theorem~\ref{thm:doubled_unitary_transport}, with $U=F$ on the
    doubled space, gives~(\ref{eq:doubled_swap_conclusions}).
\end{proof}

If $\Delta>0$ in Theorem~\ref{thm:doubled_partial_swap_gap}, then
Theorem~\ref{thm:positive_gap_uniqueness} applies to the doubled Hamiltonian
before and after the physical partial swap. In particular,
\begin{align}
    \ker(\mathbb H-2E_0I)&=\spn_{\C}\{\Psi\}, &
    \ker(\mathbb H'-2E_0I)&=\spn_{\C}\{\Phi\}.
    \label{eq:doubled_unique_ground_lines}
\end{align}
Both eigenspaces in~(\ref{eq:doubled_unique_ground_lines}) have dimension
one. Strict positivity is necessary: on $\C^2$, the Hamiltonian $H=E_0I$
satisfies the lower-gap inequality for $\Delta=0$, while its eigenspace
at $E_0$ has dimension two.
